\chapter{Hypergeometric Functions}
\label{ch:15}
\section{Krawtchouk polynomials as hypergeometric functions}
\label{sec:15-1}
Comparing coefficients in \eqref{eq:gen} gives the terminating form
\begin{equation}\label{eq:K-hypergeometric}
K_n(j;N)=\binom Nn\,{}_2F_1\!\left(-n,-j;-N;2\right),\qquad 0\le n,j\le N.
\end{equation}
The denominator parameter $-N$ causes no difficulty: the series stops at
$\min(n,j)\le N$.

\subsection{The bilinear formula}
For $t\notin\{0,-1,-2,\ldots\}$ and $A=1-(1-v)w$, Euler integrals give
\[
\sum_{n\ge0}\frac{(t)_n}{n!}w^n
{}_2F_1(-n,X;t;v){}_2F_1(-n,Y;t;v)
=(1-w)^{X+Y-t}A^{-X-Y}{}_2F_1\!\left(X,Y;t;\frac{wv^2}{A^2}\right),
\]
initially near $w=0$; it converges absolutely when
$|w|\max\{1,|1-v|^2\}<1$. Setting $X=-k$, $Y=-\ell$ and then $t=-N$
gives a finite polynomial identity. With $v=2$, $w=-r$ and
\eqref{eq:K-hypergeometric},
\begin{equation}\label{eq:closed-bilinear}
\sum_{n=0}^N\frac{r^n}{\binom Nn}K_n(k;N)K_n(\ell;N)
=(1+r)^{N-k-\ell}(1-r)^{k+\ell}\,
{}_2F_1\!\left(-k,-\ell;-N;-\frac{4r}{(1-r)^2}\right).
\end{equation}
Apparent poles cancel when $k+\ell>N$. The left side is the kernel
$m_{r,N}$ of \eqref{eq:ency-finite-mehler}; Chapter~\ref{ch:11} uses this
closed form for the Mehler limit. With $X=-x$, $Y=-y$ and $v=1-c^{-1}$,
the same identity is the Meixner bilinear kernel.

\section{Addition formulas}
\label{sec:15-2}
% TODO: Develop this section; add verified references, examples, and exercises.

\noindent\textbf{ADDITION FORMULAS NEEDS TO BE DONE.}


